\section{Uniform inputs and the exact one-switch minimax}\label{sec:one-switch}

We first solve a benchmark family for every budget with fewer activation
blocks than modes. This family gives a dimension-dependent obstruction to
small error. We then prove that its obstruction, together with an omitted-mode
obstruction, determines the full continuous one-switch minimax for at least
three modes.

\subsection{Exact uniform-input error in continuous time}

Write $A^{\mathrm{unif}}_i(t)=t/n$ and let
$E_{n,s}(T)=\OPT_s(A^{\mathrm{unif}})$.

\begin{theorem}[Uniform controls]\label{thm:uniform}
Let $n\ge2$ and $0\le s\le n-2$. With $k=s+1$ and $r=n/(n-1)$,
\begin{equation}\label{eq:uniform-value}
 E_{n,s}(T)=T\max\left\{\frac1n,\frac{1}{n(r^k-1)}\right\}.
\end{equation}
A schedule using $k$ distinct modes attains this value, although competing
schedules may repeat modes.
\end{theorem}
\begin{proof}
Consider any schedule with $m\le k<n$ positive-length blocks and error $E$.
At least one mode is unused, so $E\ge T/n$. Let $0=t_0<\cdots<t_m=T$
be its block endpoints. At the end of block $j$, the occupation of its active
mode $p_j$ is at least $t_j-t_{j-1}$, even if $p_j$ occurred earlier.
The inequality $W_{p_j}(t_j)-t_j/n\le E$ implies
\begin{equation}\label{eq:uniform-reach}
 t_j\le r t_{j-1}+rE.
\end{equation}
Iterating from zero yields
$T\le nE(r^m-1)\le nE(r^k-1)$ and proves both lower bounds.

For attainment, set
\[
 E_0=\frac{T}{n(r^k-1)},\qquad
 t_j=T\frac{r^j-1}{r^k-1}\quad(0\le j\le k),
\]
and choose a distinct mode on each block. Before its activation, the
selected mode has discrepancy $t/n$; during activation its discrepancy
decreases; after activation it increases again. At its block endpoint,
\[
 A_{p_j}(t_j)-W_{p_j}(t_j)
 =\frac{t_j}{n}-(t_j-t_{j-1})=-E_0.
\]
Every positive discrepancy is at most $T/n$, and no selected mode's
negative discrepancy exceeds $E_0$. Unused modes also have error at most
$T/n$. This proves attainment of~\eqref{eq:uniform-value}.
\end{proof}

For one switch, this gives
\begin{equation}\label{eq:uniform-one-switch}
 E_{n,1}(T)=T\max\left\{\frac1n,
 \frac{(n-1)^2}{n(2n-1)}\right\}\qquad(n\ge3).
\end{equation}
For every fixed $s\ge0$, the uniform value tends to $T/(s+1)$ as
$n\to\infty$. Indeed, with fixed $k=s+1$,
$(1+1/(n-1))^k-1=k/(n-1)+O_k(n^{-2})$, so
$n(r^k-1)\to k$. This is an obstruction supplied by one input family;
it does not presume that this family maximizes the error over all inputs.

\subsection{A scalar recurrence on uniform grids}

Let $E^{\mathrm{unit}}_{n,s}(N)$ be the grid instance optimum for the uniform
control on $N$ unit cells.

\begin{theorem}[Exact integer recurrence]\label{thm:uniform-grid}
Under $n\ge2$, $0\le s\le n-2$, and $N\ge1$, define, for integers $K\ge N$,
\begin{equation}\label{eq:uniform-grid-recurrence}
 b_0(K)=0,\qquad
 b_{j+1}(K)=\left\lfloor\frac{n b_j(K)+K}{n-1}\right\rfloor.
\end{equation}
If $K_*$ is the least such integer for which $b_{s+1}(K_*)\ge N$, then
$E^{\mathrm{unit}}_{n,s}(N)=K_*/n$. Moreover,
\begin{equation}\label{eq:uniform-grid-gap}
 E_{n,s}(N)\le E^{\mathrm{unit}}_{n,s}(N)
 <E_{n,s}(N)+\frac{n-1}{n}.
\end{equation}
\end{theorem}
\begin{proof}
At integer times, every discrepancy is a multiple of $1/n$; by
Lemma~\ref{lem:endpoint} an optimal error is therefore $K/n$ for an integer
$K$. Every admissible schedule omits a mode, giving $K\ge N$.
The block-end argument~\eqref{eq:uniform-reach}, with integer block endpoints,
gives
$t_j\le\lfloor (n t_{j-1}+K)/(n-1)\rfloor$.
Monotonicity of this recurrence implies $t_j\le b_j(K)$ and hence necessity
of $b_{s+1}(K)\ge N$. Fewer blocks can be padded by zero-length blocks;
alternatively the nondecreasing iterates give the same conclusion.

Conversely, suppose the condition holds. Put $t_j=\min\{N,b_j(K)\}$ and
stop when $N$ is first reached. Since the recurrence cannot leave zero if
$b_1(K)=0$, the assumed condition implies $b_1(K)\ge1$, or $K\ge n-1$.
It follows that the iterates strictly increase before truncation. Assign a
distinct mode to each positive-length block. Its block-end inequality gives
$W_i(t_j)-t_j/n\le K/n$. Before and after that block, its discrepancy is
nondecreasing, so this bounds its negative part throughout. Positive
errors of selected and unused modes are at most $N/n\le K/n$. There are
at most $s+1$ blocks, proving sufficiency.

The predicate is monotone in $K$, and $K=N(n-1)$ already reaches $N$ in
one step. Thus $K_*$ exists and binary search may use
$N\le K\le N(n-1)$, with $s+1$ recurrence evaluations per trial.
To prove the strict gap, write $E=E_{n,s}(N)$ and
$C=\lceil nE\rceil$, and test $K=C+n-2\ge N$. For every integer $b$,
\[
 \left\lfloor\frac{nb+C+n-2}{n-1}\right\rfloor
 =\left\lceil\frac{nb+C}{n-1}\right\rceil
 \ge r b+rE.
\]
The discrete recurrence therefore dominates the continuous recurrence
starting at zero, which reaches at least $N$ in $s+1$ steps by
\eqref{eq:uniform-value}. Consequently
$E^{\mathrm{unit}}_{n,s}(N)\le K/n<E+(n-1)/n$.
The lower inequality follows by restricting continuous switching times to
integer times for the same input.
\end{proof}

\subsection{The continuous one-switch minimax}

For two distinct modes $p,q$, let $W^{p,q,\tau}$ activate $p$ before time
$\tau$ and $q$ afterward, allowing $\tau=0$ or $T$. These endpoints include
constant schedules.

\begin{lemma}[The three-term formula]\label{lem:one-switch-error}
For $n\ge2$ and $0\le\tau\le T$, taking the maximum over omitted
terminal masses as zero when $n=2$,
\begin{equation}\label{eq:one-switch-error}
 D(A,W^{p,q,\tau})=
 \max\left\{\max_{i\notin\{p,q\}}m_i,\quad
 \tau-A_p(\tau),\quad T-m_q-\tau\right\}.
\end{equation}
\end{lemma}
\begin{proof}
The omitted modes increase to their terminal masses. The discrepancy of $p$
decreases from zero to $A_p(\tau)-\tau$ and then increases to $m_p-\tau$.
Its possible positive final error is bounded by $T-m_q-\tau$, because
$m_p+m_q\le T$. The discrepancy of $q$ increases to $A_q(\tau)$ and then
decreases to $m_q-(T-\tau)$. Its positive error at activation is at most
$\tau-A_p(\tau)$, because $A_p(\tau)+A_q(\tau)\le\tau$.
These comparisons bound every endpoint extremum by the right-hand side.
Conversely, the first two displayed terms are actual errors. If the last
term is positive it is an actual negative terminal discrepancy; if it is
nonpositive it is dominated by the nonnegative second term. This proves the
equality, including $\tau=0,T$.
\end{proof}

\begin{theorem}[Exact continuous one-switch minimax]\label{thm:one-switch}
For every $n\ge3$,
\begin{equation}\label{eq:one-switch-minimax}
 F_{n,1}(T)=H_n(T):=
 T\max\left\{\frac13,\frac{(n-1)^2}{n(2n-1)}\right\}.
\end{equation}
For $n=3,4$, three consecutive distinct pure modes, each of duration $T/3$,
attain the maximum. For $n\ge5$, the uniform control attains it. A schedule
with the asserted error bound can be constructed using terminal masses and
cumulative allocations at at most two additional times.
\end{theorem}
\begin{proof}
By scaling take $T=1$ and put $E=H_n(1)$. Directly from the formula,
$1/3\le E<1/2$.

First suppose some $m_q\ge E$. At most one other mode has mass strictly
larger than $E$: two such modes together with $q$ would have total exceeding
$3E\ge1$. Choose $p\ne q$ to include that mode if it exists, and choose any
other mode otherwise. All omitted masses are at most $E$. With
$\tau=\max\{0,1-m_q-E\}$ we have
$0\le\tau\le\max\{0,1-2E\}\le E$. Thus
$\tau-A_p(\tau)\le E$ and $1-m_q-\tau\le E$, and
Lemma~\ref{lem:one-switch-error} gives a feasible schedule.

It remains to treat the case $m_i<E$ for every mode. Let $q$ index a largest
mass and $r$ a second largest, breaking ties arbitrarily. Define
\[
 t_q=1-m_q-E,\qquad t_r=1-m_r-E,\qquad h_i=1-m_i-2E.
\]
Since $m_i<E<1/2$, the times satisfy $0<t_q\le t_r<1$.
Consider all schedules $p\to q$, $p\ne q$, at time $t_q$, and also
$q\to r$ at time $t_r$. Their omitted masses are below $E$, and their
last terms in~\eqref{eq:one-switch-error} equal $E$. If all these schedules
had error strictly larger than $E$, the middle terms would force
\[
 A_p(t_q)<h_q\quad(p\ne q),\qquad
 A_q(t_q)\le A_q(t_r)<h_r.
\]
Summing gives $t_q<(n-1)h_q+h_r$, or
\begin{equation}\label{eq:one-switch-contradiction}
 (2n-1)E<n-1-(n-2)m_q-m_r.
\end{equation}
But $m_q\ge1/n$ and $m_r\ge(1-m_q)/(n-1)$. For $n\ge3$ the coefficient
$n-2-1/(n-1)$ is positive, so
\[
 (n-2)m_q+m_r
 \ge\left(n-2-\frac1{n-1}\right)m_q+\frac1{n-1}
 \ge\frac{n-1}{n}.
\]
Equation~\eqref{eq:one-switch-contradiction} would imply
$(2n-1)E<(n-1)^2/n$, contrary to the definition of $E$. At least one
candidate therefore has error at most $E$.

For a matching lower bound of $1/3$, use three consecutive distinct pure
modes of duration $1/3$. Every schedule with at most two blocks omits at
least one of these modes and incurs its terminal mass. The uniform-input
value in~\eqref{eq:uniform-one-switch} supplies the other lower bound.
Finally,
\[
 \frac{(n-1)^2}{n(2n-1)}-\frac13
 =\frac{n^2-5n+3}{3n(2n-1)}
\]
is negative at $n=3,4$ and positive for every integer $n\ge5$, giving the
claimed transition and extremizers.
\end{proof}

For an arbitrary horizon $T>0$, set $E=H_n(T)$ and $t_i=T-m_i-E$.
The construction requires no search over switching times. Let $q$ and $r$
index a largest and second-largest terminal mass. If $m_q\ge E$, select
$r\to q$ at $\max\{0,T-m_q-E\}$. Otherwise evaluate $A(t_q)$ and
choose $p\ne q$ maximizing $A_p(t_q)$. If its candidate fails, the
$q\to r$ candidate at $t_r$ succeeds. Apart from cumulative integration,
this requires $O(n)$ arithmetic operations and comparisons. The proof covers
ties, zero terminal masses, and switching at an endpoint.

\begin{corollary}[One-switch grid upper bound]\label{cor:one-switch-grid}
For every grid $\mathcal T$ and every measurable relaxed input,
\[
 \OPT_1^{\mathcal T}(A)\le H_n(T)+\frac{\bar\Delta}{2}\qquad(n\ge3).
\]
The coefficient $H_n(T)/T$ cannot be decreased in a bound whose grid term
vanishes under refinement, uniformly over grids and inputs.
\end{corollary}
\begin{proof}
Choose the schedule from Theorem~\ref{thm:one-switch} and move its switch to
a closest grid point. The displacement is at most $\bar\Delta/2$. The
cumulative occupation of either affected mode changes by at most that
displacement, so the triangle inequality gives the bound. No new switch is
created. For sharpness of the leading coefficient, use the uniform extremizer
when $n\ge5$. When $n=3,4$, use the three pure blocks on successively refined
grids containing $T/3$ and $2T/3$. Each input belongs to its grid class, and
its grid optimum is at least its continuous optimum $H_n(T)$.
\end{proof}

\subsection{Exact three-cell minimax and a published conjecture}

\begin{proposition}[Three unit cells]\label{prop:three-cells}
For $n\ge3$, $F^{\mathrm{unit}}_{n,1}(3)=2-3/n$.
\end{proposition}
\begin{proof}
For arbitrary relaxed input choose $q$ with a largest terminal mass and $p$
with a second largest, and use the three-cell word $(p,q,q)$. Every omitted
mass is at most $1$, because the masses sum to $3$. The first-cell error is
at most $1$. At time $2$, the selected modes have occupation $1$ and
cumulative allocations in $[0,2]$, hence errors at most $1$; omitted modes
still have errors at most $1$. At time $3$, $m_p\le3/2$ and its occupation
is $1$, so its error is at most $1$. The final mode satisfies
$3/n\le m_q\le3$ and has occupation $2$, giving error at most
$\max\{1,2-3/n\}=2-3/n$. Endpoint monotonicity proves the upper bound.
For uniform input, every schedule with at most one switch gives some mode
at least two cells. Its negative terminal discrepancy is at least
$2-3/n$. Thus the upper bound is attained.
\end{proof}

Sager and Zeile~\cite[Conjecture~1, equation~(7.6)]{sager-zeile2021} propose,
for $n>2$ and $1\le s\le N-2$, the worst-case grid value
\begin{equation}\label{eq:published-conjecture}
 \begin{cases}
 T/(s+2)+\bar\Delta/2,&s\le n-2,\\
 T/(2s+4-n)+\bar\Delta/2,&s>n-2.
 \end{cases}
\end{equation}
Their initial activation is free, as in our formulation
\cite[Definition~10]{sager-zeile2021}. The following examples satisfy all
of the stated restrictions and disprove the first branch even as an upper
bound. No conclusion about every case of its second branch is needed here.

\begin{example}[Five modes and one switch]\label{ex:five-mode-conjecture}
Take $n=5$, $T=N=45$ unit cells, and $\alpha_i=1/5$ throughout.
For any distinct pair of modes and an integer switch time $\tau$, the exact
error from~\eqref{eq:one-switch-error} is
\[
 D(\tau)=\max\{9,4\tau/5,36-\tau\}.
\]
The two nonconstant terms intersect at $\tau=20$, with value $16$.
For $\tau\le20$ the last term is at least $16$; for $\tau\ge20$ the
middle term is at least $16$. Constant schedules are included by the endpoint
choices. Thus the grid optimum is $16$, whereas
\eqref{eq:published-conjecture} gives $45/3+1/2=31/2$.
On $45m$ unit cells the same construction has optimum $16m$, exceeding the
conjectured $15m+1/2$ for every integer $m\ge1$. Equivalently, on a fixed
horizon $T$ with $45m$ equal cells its optimum is $16T/45$, while the
conjectured bound tends to $T/3$. The failure persists under refinement.
\end{example}

\begin{example}[Smallest admissible number of cells]\label{ex:three-cell-conjecture}
On three unit cells and with seven uniform modes, the exact grid error is
$11/7$ by Proposition~\ref{prop:three-cells}. It exceeds the conjectured
$3/3+1/2=3/2$. The restriction $1\le s\le N-2$ excludes grids with fewer
than three cells, so this witness uses the smallest admissible number of
cells.
\end{example}

More generally, for every fixed $s\ge1$, Theorem~\ref{thm:uniform} approaches
$T/(s+1)>T/(s+2)$ as the mode count grows. For sufficiently many modes, its
continuous uniform value already exceeds the first branch's proposed leading
term. A sufficiently fine uniform grid then makes $\bar\Delta/2$ smaller
than this strict excess; restricting the uniform input to grid schedules
cannot reduce its optimum. Thus this obstruction occurs for every fixed
positive switch budget, not only for one switch. In the one-switch case,
Theorem~\ref{thm:one-switch} and Corollary~\ref{cor:one-switch-grid} give an
exact continuous replacement and a valid grid upper bound. The latter is
not asserted to be an exact finite-grid equality.
